% ==============================================================================
% ANTECEDENTES
% ==============================================================================
% Contenido redactado y provisto directamente por el estudiante (no adaptado
% del protocolo). Se aplicaron dos correcciones editoriales antes de
% incorporarlo:
%   1. La segunda tabla (comparativa de soluciones RASP) usaba columnas "l"
%      sin ancho fijo con encabezados largos ("Detección de variantes
%      desconocidas"), lo que se sale del margen de la página. Se cambió a
%      columnas p{} con \small, igual que la primera tabla.
%   2. La entrada bibliográfica mereani2019 traía \type{@article} con campos
%      de conferencia (booktitle, organization); se corrigió a
%      \type{@inproceedings}, que es el tipo correcto en biblatex para esos
%      campos.
\ifdefined\modoensamble
\else
	\input{0-preambulo_standalone}
\fi

\chapter{Antecedentes}

\section{Estado del arte del problema}
La detección automatizada de ataques a aplicaciones web mediante técnicas de aprendizaje
automático ha sido un área de investigación activa desde al menos la década pasada, pero su
consolidación como línea aplicable a arquitecturas de producción reales es reciente. La literatura
puede organizarse en tres corrientes que convergen directamente en el problema que aborda este
trabajo: (i) modelos de aprendizaje automático especializados en la detección de un vector de
ataque específico (inyección SQL, XSS, path traversal o inyección de comandos) evaluados sobre
conjuntos de datos aislados; (ii) soluciones comerciales de protección en tiempo de ejecución
(RASP) orientadas a entornos empresariales con presupuesto de seguridad dedicado; y (iii) el
contexto regulatorio y de madurez tecnológica de Guatemala, donde tanto la oferta de herramientas
de seguridad accesibles como el marco legal que las incentivaría se encuentran en un estado
incipiente. Ninguna de estas tres corrientes, por separado, resuelve el problema identificado en la
sección 2 del protocolo de este trabajo: la ausencia de un mecanismo de detección multivector,
integrado nativamente en el runtime de Node.js, distribuible sin costo y operable por equipos sin
un analista de seguridad dedicado. Este capítulo revisa los antecedentes más relevantes de cada
corriente, identifica sus aportes y limitaciones, y sintetiza el vacío de investigación que justifica la
propuesta de \textit{logSguarDian}.

\section{Contexto de ciberseguridad y marco regulatorio en Guatemala}
El ecosistema tecnológico guatemalteco presenta una combinación de crecimiento acelerado en la
adopción de aplicaciones web y ausencia casi total de controles de seguridad estructurados. El
Índice de Delincuencia Organizada reporta un aumento de los delitos cibernéticos en el país que atribuye a la falta de regulación y
control en internet \cite{lahora2025ley6347}. Esta carencia regulatoria no es un dato
anecdótico: constituye en sí misma un antecedente directo del problema, porque en ausencia de
un marco legal que obligue a las organizaciones a invertir en controles mínimos, la adopción de
herramientas de seguridad depende exclusivamente de su accesibilidad económica y operativa, lo
que es precisamente el vacío que este proyecto busca cerrar.

Guatemala cuenta desde 2024 con la Iniciativa de Ley 6347, "Ley de Ciberseguridad", que propone
la tipificación de nuevos delitos informáticos y la creación de un Centro de Respuesta a Incidentes
de Seguridad Informática (CSIRT-GT) como órgano técnico nacional \cite{congreso2025dictamen}.
La iniciativa recibió dictamen favorable de la Comisión de Asuntos de Seguridad Nacional en
agosto de 2025 y fue aprobada en segunda lectura por el pleno del Congreso el 23 de septiembre
de 2025; sin embargo, a la fecha de redacción de este documento permanece sin aprobación final,
con su tercer debate reactivado hasta abril de 2026 tras más de seis meses de estancamiento
\cite{lahora2026estancada}. Un análisis técnico-jurídico de la Alianza por un Congreso Eficiente sobre la iniciativa 6347 señaló
además ambigüedades en varios tipos penales que comprometen el principio de legalidad y recomendó
una revisión legislativa antes de su aprobación \cite{cien2025boletin}. Este antecedente
regulatorio es relevante para el posicionamiento del proyecto por dos razones. Primero, confirma
que, incluso si la ley se aprueba, su enfoque es punitivo y de respuesta institucional a incidentes,
no preventivo a nivel de aplicación; no existe en la iniciativa ningún mecanismo que reduzca la
brecha de herramientas de defensa activa para las startups que constituyen la mayoría del tejido
empresarial guatemalteco. Segundo, mientras el marco legal no entre en vigor, la adopción de
controles de seguridad en el país seguirá siendo enteramente voluntaria, lo que refuerza la
necesidad de que dichos controles sean gratuitos y de instalación trivial para que exista algún
incentivo real de adopción.

\section{Trabajos académicos en detección de ataques web mediante aprendizaje automático}
La aplicación de aprendizaje automático a la detección de ataques específicos de aplicaciones web
ha producido resultados cuantitativos sólidos, pero con limitaciones estructurales consistentes en
su diseño experimental. El Cuadro \ref{cuadro:trabajosprevios} resume los trabajos más
relevantes para el contexto de este proyecto.

\begin{table}[h]
\small
\begin{tabular}{|p{2.1cm}|p{1.6cm}|p{2.2cm}|p{1.6cm}|p{1.6cm}|p{3.3cm}|}
\hline
\textbf{Trabajo} & \textbf{Vector} & \textbf{Técnica} & \textbf{Dataset} & \textbf{Métrica} & \textbf{Limitación principal} \\ \hline
Peralta-Garcia et al. \cite{peralta2024sql} & SQL Injection & Random Forest, árbol de decisión, SVM & Dataset público de 22{,}764 registros & Accuracy = 99~\%; F1 = 98~\% (Random Forest) & Solo cubre SQLi; el modelo se despliega en una arquitectura de microservicios que valida peticiones, es decir, como servicio separado de la aplicación \\ \hline
Mereani \& Howe \cite{mereani2019} & XSS & SVM, k-NN, Random Forest & Código JavaScript etiquetado & No reportado & Clasifica código del lado del cliente, no parámetros de peticiones entrantes al servidor \\ \hline
Wang et al. \cite{seyyar2024} & Command Injection & CCBA: CNN de doble canal, BiLSTM y atención & Dataset propio de inyección de comandos; evaluación cruzada en CSIC 2010 \cite{torrano2010csic} y HTTPParam & Accuracy = 99.3~\%; recall = 98.2~\% (dataset propio) & Modelo de aprendizaje profundo evaluado en una GPU; no reporta latencia de inferencia ni despliegue como middleware \\ \hline
SAP node-rasp \cite{saparchive2019} & SQLi, NoSQLi, Path Traversal & Análisis léxico con reglas (no ML) & No aplica & No reportado & Requiere modificar el runtime de Node.js; abandonado desde LTS 10; no distribuible como paquete npm \\ \hline
Durmuşkaya y Bayraklı \cite{durmuskaya2025waf} & SQLi, XSS, CMDi, LFI & Árbol de decisión (comparado con KNN, regresión logística, Naive Bayes y SVM) & CSIC 2010 \cite{torrano2010csic}, HTTPParams 2015 y conjunto propio & Accuracy = 93.27~\% (WAF en vivo) & Opera como proxy inverso externo (mitmproxy, Python); evaluado contra DVWA; sin cifra de latencia reportada \\ \hline
\end{tabular}
\caption[Comparativa de trabajos previos en detección de ataques web con ML]{Comparativa de trabajos previos en detección de ataques a aplicaciones web mediante aprendizaje automático}
\label{cuadro:trabajosprevios}
\end{table}
{\footnotesize Nota. Elaboración propia con base en los trabajos citados en cada fila. Las métricas son las que reporta cada autor sobre sus propios datos y no son directamente comparables entre sí. WAF: \textit{Web Application Firewall}; SQLi: \textit{SQL Injection}; XSS: \textit{Cross-Site Scripting}; CMDi: \textit{Command Injection}; LFI: \textit{Local File Inclusion}.}

El análisis de estos cinco trabajos revela un patrón consistente antes que hallazgos aislados. Con la excepción parcial de Durmuşkaya y Bayraklı \cite{durmuskaya2025waf}, quienes despliegan en su
WAF un único árbol de decisión binario aunque evalúan el desempeño por vector con modelos separados,
la totalidad de los trabajos revisados entrena y evalúa un clasificador por vector de ataque de forma
independiente, sin considerar que en un escenario de producción real un mismo endpoint puede
recibir simultáneamente intentos de inyección SQL, XSS y path traversal. Incluso el trabajo más cercano en arquitectura a la propuesta de este proyecto, el WAF de Durmuşkaya
y Bayraklı \cite{durmuskaya2025waf}, resuelve el problema mediante un proxy inverso externo
(mitmproxy, implementado en Python), lo que introduce un
salto de red adicional y no captura el contexto interno del proceso Node.js (memoria del \texttt{event
loop}, tiempos de procesamiento internos, encadenamiento de middlewares) que sí está disponible
para una librería integrada directamente como middleware. Ninguno de los trabajos revisados
reporta el impacto de su solución sobre la latencia del servidor protegido bajo carga sostenida, dato
que es un criterio de aceptación explícito de este proyecto.

Es pertinente señalar también una limitación metodológica: las métricas reportadas no son homogéneas entre los trabajos revisados (\textit{accuracy} y F1 en
\cite{peralta2024sql}; \textit{accuracy} y \textit{recall} en \cite{seyyar2024}; \textit{accuracy} en \cite{durmuskaya2025waf}), lo que impide una comparación directa, y confiar en la
\textit{accuracy} sola es una práctica problemática cuando los datasets de ataques suelen
estar desbalanceados a favor del tráfico legítimo, ya que una accuracy alta puede ocultar una tasa
de falsos negativos inaceptable para un sistema de seguridad en producción.

\section{Soluciones RASP existentes: antecedentes tecnológicos}
\label{sec:rasp-antecedentes}
\label{sec:rasp-existentes}
Además de los antecedentes académicos, es necesario revisar los desarrollos tecnológicos
disponibles comercialmente bajo la categoría de Runtime Application Self-Protection (RASP), ya
que representan el estado práctico de la industria frente al mismo problema. El Cuadro
\ref{cuadro:rasp} compara las soluciones más relevantes.

\begin{table}[h]
\small
\begin{tabular}{|p{2.4cm}|p{1.6cm}|p{2cm}|p{1.8cm}|p{2.6cm}|p{1.8cm}|}
\hline
\textbf{Solución} & \textbf{Soporte Node.js} & \textbf{Costo} & \textbf{Plataforma externa} & \textbf{Detección de variantes desconocidas} & \textbf{Operación autónoma} \\ \hline
Contrast Protect \cite{contrastsecurity} & Sí & Licencia empresarial & Sí & Sí (IAST + RASP) & No \\ \hline
Datadog App and API Protection (antes ASM; incorpora Sqreen) \cite{datadog2026nodejs,datadog2021sqreen} & Sí & Suscripción & Sí & Parcial & No \\ \hline
OpenRASP (Baidu) \cite{baidu2026openrasp} & No & Gratuito & No & Limitado & Sí \\ \hline
\textit{logSguarDian} (propuesto) & Sí & Gratuito (npm) & No & Sí (medida sobre corpus limpio) & Sí \\ \hline
\end{tabular}
\caption[Comparativa de soluciones RASP existentes]{Comparativa entre soluciones RASP comerciales existentes y la propuesta de este trabajo}
\label{cuadro:rasp}
\end{table}
{\footnotesize Nota. Elaboración propia con base en la documentación de cada producto
\cite{contrastsecurity,datadog2026nodejs,baidu2026openrasp}. Las columnas de costo, plataforma
externa y detección de variantes desconocidas reflejan el criterio del autor y no una tarifa
publicada.}

Esta comparación establece un antecedente tecnológico claro: existe tecnología RASP madura en
el mercado, pero ninguna de las soluciones revisadas combina soporte nativo para Node.js,
gratuidad, y operación sin dependencia de una plataforma externa de análisis. Contrast Protect y Datadog, las opciones comerciales revisadas con soporte para Node.js, dependen, según el criterio del autor, de una plataforma del proveedor, lo que introduce costo recurrente y consideraciones de
privacidad de datos que resultan prohibitivas para una startup guatemalteca sin presupuesto de
seguridad. OpenRASP, la única alternativa gratuita entre las soluciones revisadas, ofrece agentes solo para
Java y PHP \cite{baidu2026openrasp}, y su enfoque
de detección basado en reglas es más limitado frente a variantes desconocidas que un clasificador
entrenado sobre la distribución estadística de los ataques.

\section{Análisis comparativo y vacíos identificados}
La revisión conjunta de los antecedentes académicos y tecnológicos permite identificar tres brechas
consistentes que ningún trabajo previo resuelve de forma simultánea.

La primera es la \textbf{brecha de integración en runtime}. Los trabajos académicos de detección de
SQLi y XSS revisados operan sobre datasets de cadenas de texto o payloads aislados, sin
integrarse en el flujo de ejecución del servidor; incluso el WAF de Durmuşkaya y Bayraklı \cite{durmuskaya2025waf}, la propuesta arquitectónicamente
más cercana entre las revisadas, opera como intermediario de red externo y no como componente interno del proceso Node.js. Esto significa que ninguno de estos
sistemas ha sido validado bajo restricciones de latencia de producción con acceso al contexto
interno completo de una petición HTTP.

La segunda es la \textbf{brecha de cobertura multivector con enfoque unificado dentro del propio
proceso de la aplicación}. Aunque Durmuşkaya y Bayraklı \cite{durmuskaya2025waf} sí cubren varios vectores de ataque (SQLi,
XSS, inyección de comandos y LFI) en un mismo sistema, lo hacen desde un proxy inverso externo,
entrenado con CSIC 2010 \cite{torrano2010csic}, HTTPParams 2015 y un conjunto propio, y evaluado contra DVWA, una
aplicación de formularios, y no contra APIs REST modernas que transmiten payloads en cuerpos JSON en
lugar de parámetros de URL. El resto de
los trabajos académicos revisados aborda un único vector por sistema.

La tercera es la \textbf{brecha de viabilidad operacional para organizaciones sin equipo de
seguridad}, que es simultáneamente un vacío académico y comercial. Los trabajos académicos
revisados reportan métricas de rendimiento del modelo pero no reportan impacto en latencia del
servidor anfitrión ni consumo de memoria del proceso; las soluciones comerciales que sí resuelven
estos aspectos operacionalmente lo hacen a un costo de licencia y de dependencia de
infraestructura externa incompatible con el perfil de una startup guatemalteca. SAP node-rasp
\cite{saparchive2019} es el único antecedente que intentó específicamente la protección en
runtime de Node.js sin dependencias externas, pero su enfoque de modificar el binario del runtime
lo hizo incompatible con entornos de producción estándar y fue abandonado sin soporte desde
versiones posteriores a Node.js LTS 10.

\section{Síntesis y relación con la propuesta de investigación}
Los antecedentes revisados confirman que el problema abordado por este trabajo no es la ausencia
de investigación en detección de ataques web mediante aprendizaje automático, tema que cuenta
con una literatura considerable y en expansión activa, como lo demuestra la publicación de Durmuşkaya y Bayraklı \cite{durmuskaya2025waf} en 2025. El vacío real
es la ausencia de una solución que combine simultáneamente cuatro condiciones que hasta ahora
solo han sido resueltas de forma parcial y nunca en conjunto: integración nativa en el runtime de
la aplicación en lugar de un componente de red externo, cobertura de múltiples vectores de ataque
en un único pipeline de detección, gratuidad y distribución mediante un mecanismo trivial de
instalación (\texttt{npm install}), y evaluación explícita del impacto en latencia y consumo de recursos
como criterio de aceptación del sistema. A esto se suma el vacío de contexto local: en ausencia de
un marco regulatorio vigente en Guatemala \cite{lahora2026estancada} y de herramientas
adaptadas a la realidad operativa de las startups tecnológicas del país, la adopción de controles de
seguridad depende enteramente de que existan opciones accesibles, lo cual no ha sido el caso
hasta ahora. \textit{logSguarDian} se posiciona precisamente en la intersección de estos vacíos: no
compite con las soluciones RASP empresariales ni pretende sustituir a los WAF perimetrales
existentes, sino que ocupa el espacio que ninguno de los antecedentes revisados cubre
completamente, proporcionando un mecanismo de detección multivector, integrado, gratuito y
operable sin equipo de seguridad dedicado.

\ifdefined\modoensamble
\else
	\section*{Referencias}
	\printbibliography[heading=none]
	\end{document}
\fi
